\documentclass[11pt,letterpaper]{article}
\usepackage[T1]{fontenc}
\usepackage{lmodern,amsmath,amssymb,amsthm,mathtools,mathrsfs}
\usepackage[margin=0.94in,headheight=15pt]{geometry}
\usepackage{microtype,booktabs,longtable,array,enumitem,fancyhdr,xcolor,xurl}
\usepackage[colorlinks=true,linkcolor=blue!42!black,citecolor=blue!42!black,urlcolor=blue!42!black]{hyperref}
\usepackage{bookmark}
\numberwithin{equation}{section}
\newtheorem{theorem}{Theorem}[section]
\newtheorem{lemma}[theorem]{Lemma}
\newtheorem{proposition}[theorem]{Proposition}
\newtheorem{corollary}[theorem]{Corollary}
\theoremstyle{definition}\newtheorem{definition}[theorem]{Definition}
\theoremstyle{remark}\newtheorem{remark}[theorem]{Remark}
\newtheorem{example}[theorem]{Example}
\newcommand{\Q}{\mathbb Q}\newcommand{\C}{\mathbb C}\newcommand{\N}{\mathbb N}
\newcommand{\D}{\mathcal D}\newcommand{\NF}{\mathcal N}\newcommand{\Prim}{\mathcal T}
\newcommand{\A}{\mathcal A}\newcommand{\HH}{\mathcal H}\newcommand{\BB}{\mathcal B}
\newcommand{\dd}{\,\mathrm d}\newcommand{\ii}{\mathrm i}
\newcommand{\piy}{\partial_{y^2}}
\DeclareMathOperator{\Li}{Li}\DeclareMathOperator{\Shi}{Shi}
\DeclareMathOperator{\FP}{FP}\DeclareMathOperator{\Res}{Res}\DeclareMathOperator{\CT}{CT}
\newcommand{\src}[1]{\nolinkurl{#1}}
\newcommand{\pin}{92ee3aaf3087c0e0c29b9ed952e777de5d6accd0}
\setlength{\emergencystretch}{3em}
\setlength{\parskip}{3pt}
\setlist{itemsep=5pt,topsep=5pt}
\allowdisplaybreaks[2]
\setcounter{tocdepth}{2}
\pagestyle{fancy}\fancyhf{}
\fancyhead[L]{\small Affine Gamma polarization}
\fancyhead[R]{\small ProveIt research continuation}
\fancyfoot[C]{\thepage}
\hypersetup{pdftitle={Affine Polarization of Gamma Normal Forms},pdfauthor={Research draft prepared for Vladimir Reshetnikov with OpenAI assistance},pdfsubject={Exact Gamma coefficients, Hurwitz and Stieltjes averages, polylogarithmic kernels}}
\begin{document}
\hypersetup{pageanchor=false}
\begin{titlepage}
\centering
\vspace*{0.7cm}
{\large\sffamily PROVEIT / ANALYSIS / POLYLOGARITHMS\par}
\vspace{1.1cm}
{\Huge\bfseries Affine Polarization\\[0.15em] of Gamma Normal Forms\par}
\vspace{0.7cm}
{\Large All derivative-square coefficients,\\[0.2em]
polylogarithmic kernels, and Stieltjes averages\par}
\vspace{0.8cm}
{\large Research continuation prepared for Vladimir Reshetnikov\par}
\vspace{0.25cm}
{Mathematical derivations and computational checks with OpenAI assistance\par}
\vspace{0.25cm}
{11 October 2026\par}
\vspace{0.7cm}
\begin{minipage}{0.93\textwidth}
\small
\textbf{Abstract.}
The canonical integration-by-parts normal form of
$(-1)^m\Gamma(x)(\Gamma(x)^{-1})^{(m)}$ contains a tractable infinite
sector that is not apparent from its Bell-polynomial expansion.
We determine every coefficient of $u_0^a u_1^k u_r^2$, simultaneously
for all $r\ge2$ and all $a,k\ge0$. An affine perturbation turns the
quadratic reduction into one symmetric symbol. Its coefficients are
entire hypergeometric functions, with an elementary beta-integral
representation and a two-variable generating function. A general
symmetric-profile theorem extends the construction to products of
centered Gamma ratios.

The same kernel gives a normalized Mellin--Hurwitz transform with only
$r+1$ endpoint pairs. Its Laurent expansion yields finite formulas for
every beta-weighted generalized Stieltjes constant, with explicit
complete-harmonic coefficients and anchored primitives of every order.
A Lerch deformation gives finite polylogarithmic endpoint formulas.
We evaluate six quartic-weight polygamma integrals and prove a
Bernoulli selector basis characterized by single-zeta closure over rational
functions of the complex order. Its trigamma specialization contains one
odd zeta value plus an explicit rational number. All displayed integrals in these analytic
applications are ordinary convergent integrals.

The normal-form decomposition, the pure-square coefficients, and the
$r=2$ family are credited to prior work. The new coefficient theorem
answers a defined subcase of the predecessor's next-Gamma-projection
question; it does not solve its full mixed $u_2,u_3$ sector or the general
harmonic-tail rank conjecture. Proofs, 2,187 exact assertions, 56 numerical
diagnostics, and additive integration instructions accompany the article.
Historical priority beyond the inspected sources is not asserted.
\end{minipage}
\vfill
{\small Repository metadata anchor: \texttt{\pin}\par}
{\small Source identity and inspection limits are recorded in Section~\ref{ag:sec:audit}.\par}
\end{titlepage}
\hypersetup{pageanchor=true}
\pagenumbering{roman}
{\small\setlength{\parskip}{0pt}\tableofcontents}
\clearpage\pagenumbering{arabic}

% BEGIN sections/01_problem.tex
\section{The continuation problem and the main result}
\label{ag:sec:problem}

The manuscript's Gamma hierarchy has two distinct kinds of questions.
The first is algebraic: which differential monomials survive a specified
integration-by-parts reduction, and with which rational coefficients?
The second is analytic: how do their integrals, spectral derivatives,
and normalized primitives reduce to special functions? Neither question
should be replaced by an assertion of independence of numerical constants.

The immediate source for this continuation is \emph{Exact Harmonic-Tail
Reconstruction and Variational Gamma Identities} \cite{HTG}. Its
Section~4 determines the complete normal projection onto $u_0,u_1,u_2$.
Its final research question~5 asks for the next projection, including
$u_3^2$ with arbitrary powers of $u_1,u_2$. We solve the subfamily with
no additional $u_2$ factor, and simultaneously solve its analogue for
every derivative order $r\ge2$. In particular, the result is not only
a table at the next weight.

Let $\D u_j=u_{j+1}$ in $\Q[u_0,u_1,\ldots]$. Write $\NF$ for the
canonical projection onto monomials which are either powers of $u_0$
or have their largest positive index repeated at least twice. Define
\begin{equation}\label{ag:eq:hierarchy-intro}
 A_0=1,\qquad A_{m+1}=u_0A_m-\D A_m,\qquad N_m=\NF(A_m).
\end{equation}
The decomposition and all conventions are proved or specified in the
next section. Under $u_j=\psi^{(j)}(x)$,
\[
 A_m(x)=(-1)^m\Gamma(x)\frac{\mathrm d^m}{\mathrm dx^m}\frac1{\Gamma(x)}.
\]

\begin{theorem}[All derivative-square coefficients]\label{ag:thm:main}
Let $r\ge2$ and $a,k\ge0$, and put $m=a+2k+2r+2$. Then
\begin{equation}\label{ag:eq:coefficient}
 \boxed{
 [u_0^a u_1^k u_r^2]N_m=
 \begin{cases}
 \displaystyle
 \frac{(-1)^r m!}
 {a!\,2^{k+r+1}k!(r+1)!\prod_{j=0}^{r}(k+2j+1)},&k\text{ even},\\[7pt]
 0,&k\text{ odd}.
 \end{cases}}
\end{equation}
Equivalently, with $y=u_1t^2/2$,
\begin{equation}\label{ag:eq:normal-sector}
 \sum_{m\ge0}N_m\frac{t^m}{m!}
 =e^{tu_0}\left\{f_0(y)+\sum_{r\ge2}u_r^2t^{2r+2}g_r(y)
                          +R(t;u_1,u_2,\ldots)\right\},
\end{equation}
where every monomial of $R$ contains at least three factors from
$\{u_2,u_3,\ldots\}$, and
\begin{align}
 f_0(y)&=\frac12\left(\cosh y+\frac{\sinh y}{y}-y\Shi(y)\right),
       \label{ag:eq:f0-intro}\\
 g_r(y)&=\frac{(-1)^r}{(r+1)!}\piy^{\,r+1}\cosh y
 =\frac{(-1)^r}{(2r+2)!}\,{}_0F_1\left(;r+\frac32;\frac{y^2}{4}\right).
       \label{ag:eq:gr-intro}
\end{align}
Here $\Shi(y)=\int_0^y\sinh v\dd v/v$, and
$\piy=(2y)^{-1}\partial_y$ acts on even power series. All apparent
singularities at $y=0$ are removable.
\end{theorem}

The phrase ``quadratic sector'' refers to the number of factors having
index at least two. It does not mean the total factor degree, and it does
not discard arbitrary powers of $u_0$ or $u_1$. Formula
\eqref{ag:eq:normal-sector} is an identity of formal series, not an
approximation to a Gamma ratio on an interval. Its proof is in
Section~\ref{ag:sec:polarization}.

\subsection{The overlap and the advance}

At $k=0$, \eqref{ag:eq:coefficient} is
$(-1)^r\binom{m}{2r+2}$; the inspected predecessors already contain
these pure-square coefficients. At $r=2$ it agrees with the complete
$u_2^2$ family in \cite[Section~4]{HTG}. Those cases are checks and
inherited specializations, not additional discoveries claimed here.
The advance is the single all-$r$, all-$k$ formula and its affine-symbol
proof, together with the symmetric-profile and analytic consequences.

For example, the coefficients of $u_1^{2\ell}u_r^2$ at their corresponding
weights $4\ell+2r+2$ begin as follows. These are not coefficients at
one fixed weight.
\begin{center}
\begin{tabular}{@{}rrrrr@{}}
\toprule
$r$ & $\ell=0$ & $\ell=1$ & $\ell=2$ & $\ell=3$\\
\midrule
2&1&90&15015&4176900\\
3&$-1$&$-165$&$-40950$&$-15261750$\\
4&1&273&96390&47006190\\
5&$-1$&$-420$&$-203490$&$-127193220$\\
6&1&612&395010&310810500\\
\bottomrule
\end{tabular}
\end{center}
Multiplication by $u_0^a$ multiplies the exponential-series coefficient
by $1/a!$, as displayed in the theorem.

The analytic part is connected by the beta representation of $g_r$,
not by an unproved interchange of the entire Gamma normal-form series
with a singular integral. It yields finite Hurwitz and Lerch endpoint
formulas, every generalized Stieltjes order, and ordinary integral
identities such as
\begin{equation}\label{ag:eq:trigamma-preview}
 \int_1^3(x-1)^2(3-x)^2\psi'(x)\dd x=\frac{6\zeta(3)}{\pi^2}.
\end{equation}
Individual polynomially weighted special-function integrals can have
classical antecedents; our claim is the proved structural derivation
and its explicitly scoped extension, not universal historical priority.

% END sections/01_problem.tex


% BEGIN sections/02_normal.tex
\section{Canonical reduction, centering, and endpoint conventions}
\label{ag:sec:normal}

\subsection{The free differential algebra}
Let $\mathscr R=\Q[u_0,u_1,u_2,\ldots]$, where every polynomial uses
finitely many variables. The weight of $u_j$ is $j+1$; its factor degree
is one. The derivation $\D$ increases weight by one and preserves factor
degree. A monomial is \emph{functional} when it is a power of $u_0$,
including $1$, or when its largest positive index occurs at least twice.
Let $\mathscr F$ be the rational span of these monomials.

The following is the normal-form input used in the repository
\cite{Manuscript,HTG}. The terminology and the underlying complement
also occur in the classical differential-algebra construction
\cite[Section~4.4.1]{GRR}. We include a proof to make our coefficient
conventions independent of a particular implementation.

\begin{lemma}[The functional-monomial decomposition]\label{ag:lem:normal}
Every $P\in\mathscr R$ has a unique decomposition
\begin{equation}\label{ag:eq:normal-decomp}
 P=\NF(P)+\D\Prim(P),\qquad \NF(P)\in\mathscr F,
                  \qquad [1]\Prim(P)=0.
\end{equation}
Both maps are rational linear and preserve factor degree. On positive
weight $w$, $\NF(P)$ has weight $w$ and $\Prim(P)$ has weight $w-1$.
\end{lemma}
\begin{proof}
Suppose a monomial has its largest index $r\ge1$ exactly once. Write it
as $u_r u_{r-1}^aV$, where $V$ uses indices below $r-1$. Then
\begin{equation}\label{ag:eq:ibp-step}
 u_r u_{r-1}^aV
 =\D\left(\frac{u_{r-1}^{a+1}V}{a+1}\right)
       -\frac{u_{r-1}^{a+1}\D V}{a+1}.
\end{equation}
Every monomial of the remainder has largest index below $r$. Repeating
this replacement terminates and produces a functional remainder.
Constants of primitives can be set to zero.

For uniqueness, if a nonconstant polynomial $H$ has largest variable
index $j$, the terms $u_{j+1}\partial H/\partial u_j$ in $\D H$ are
nonzero and have largest index $j+1$ exactly once. All other terms have
smaller largest index. They cannot be canceled by a functional
polynomial. Thus $\D H\in\mathscr F$ forces $H$ constant, proving
both directness and uniqueness after the stated normalization.
\end{proof}

This is a vector-space decomposition. For instance,
$\NF(u_1)=0$ but $\NF(u_1^2)=u_1^2$; replacing $\NF(PQ)$ by
$\NF(P)\NF(Q)$ is incorrect. Another useful check is
$\NF(u_0u_4)=u_2^2$. Truncating away $u_4$ before reduction would
lose that term. These are safeguards for the present derivation, not
newly discovered errors in sources which already state the distinction.

\subsection{The centered Gamma representative}
The exponential generating function of \eqref{ag:eq:hierarchy-intro} is
\begin{equation}\label{ag:eq:Bell}
 \A(t)=\sum_{m\ge0}A_m\frac{t^m}{m!}
 =\exp\left(\sum_{j\ge0}\frac{(-1)^j u_jt^{j+1}}{(j+1)!}\right).
\end{equation}
It follows either by the exponential Bell formula or by differentiating
with respect to $t$ and using the defining recurrence. Under the digamma
substitution it is $\Gamma(x)/\Gamma(x-t)$.

Define
\begin{equation}\label{ag:eq:centered}
 \mathcal C(t)=e^{t\D/2}\A(t)
 =\exp\left(\sum_{j\ge0}
       \frac{u_{2j}t^{2j+1}}{2^{2j}(2j+1)!}\right).
\end{equation}
For any formal differential series $P(t)$,
\[
 (e^{t\D/2}-1)P(t)
 =\D\left(\sum_{n\ge1}\frac{t^n\D^{n-1}P(t)}{2^n n!}\right).
\]
Therefore $\NF\mathcal C(t)=\NF\A(t)$. Analytically the centered
representative is $\Gamma(x+t/2)/\Gamma(x-t/2)$, initially near $t=0$
away from Gamma poles.

Translation $u_0\mapsto u_0+c$ commutes with $\D$ and preserves the
functional space. Uniqueness in Lemma~\ref{ag:lem:normal} implies that
it commutes with $\NF$. Thus $\NF\A(t)$ has an Appell factor $e^{tu_0}$.
The centered series also shows factor parity: each exponent term has
odd weight and factor degree one, so every surviving monomial has
factor degree congruent to its weight modulo two.

\begin{lemma}[Appell covariance of the primitive]\label{ag:lem:appell-primitive}
Suppose $P(t)$ and $\NF P(t)$ both satisfy
$\partial_{u_0}P=tP$. If $P-\NF P=\D T$ with $[1]T=0$, then
\[
 \partial_{u_0}T-tT=c(t)
\]
for a scalar series $c(t)$. Consequently, modulo a scalar series,
$T=e^{tu_0}B(t;u_1,u_2,\ldots)$, allowing a harmless Laurent power
of $t$ if necessary.
\end{lemma}
\begin{proof}
The operators $\partial_{u_0}$ and $\D$ commute. Apply
$\partial_{u_0}-t$ to $\D T=P-\NF P$. Its result is zero, so the
kernel statement in Lemma~\ref{ag:lem:normal} makes
$\partial_{u_0}T-tT$ scalar. Solving this elementary equation in $u_0$
gives $T=e^{tu_0}B-c(t)/t$.
\end{proof}
This small observation matters in the affine-symbol proof: after
perturbation, the quadratic part of the primitive has the same
exponential weight as the density and coefficients independent of the
affine coordinate. We do not enlarge the quotient by arbitrary
$x$-dependent primitives.

\subsection{What an analytic substitution does and does not imply}
For every fixed $m$, let $T_m=\Prim(A_m)$. On an interval
$[\alpha,\beta]\subset(0,\infty)$,
\begin{equation}\label{ag:eq:ordinary-gamma-integral}
 \int_\alpha^\beta(-1)^m\Gamma(x)(\Gamma(x)^{-1})^{(m)}\dd x
 =T_m(\psi(\beta),\ldots)-T_m(\psi(\alpha),\ldots)
       +\int_\alpha^\beta N_m(\psi(x),\ldots)\dd x.
\end{equation}
If a fixed-coordinate finite part at zero is used instead, the endpoint
is $T_m(\psi(1),\ldots)-[x^0]T_m(\psi(x),\ldots)$, not zero.
This inherited endpoint convention \cite{Manuscript} is retained.

At positive integral endpoints the entries are generalized harmonic
numbers:
\begin{equation}\label{ag:eq:harmonic-substitution}
 \psi(N+1)=H_N-\gamma,\qquad
 \psi^{(j)}(N+1)=(-1)^{j+1}j!\bigl(\zeta(j+1)-H_N^{(j+1)}\bigr)
 \quad(j\ge1).
\end{equation}
These follow from the standard polygamma recurrence \cite{DLMFpsi}.
Thus the exact normal-form certificates also give finite identities
between Gamma derivatives and harmonic-tail polynomials. They do not
by themselves evaluate every integral of a remaining normal monomial.

% END sections/02_normal.tex


% BEGIN sections/03_polarization.tex
\section{Affine polarization and the complete square sector}
\label{ag:sec:polarization}

\subsection{A constant-coefficient quadratic symbol}
Let $p,q,t$ be formal parameters and make the differential substitution
\begin{equation}\label{ag:eq:affine}
 u_0=q+px+\varepsilon v(x),\qquad
 u_1=p+\varepsilon v'(x),\qquad
 u_j=\varepsilon v^{(j)}(x)\quad(j\ge2).
\end{equation}
Put $W=e^{t(q+px)}$ and $\kappa=tp$, so $W'=\kappa W$.
For a quadratic differential polynomial in $v$ with constant
coefficients, polarization introduces two derivative variables
$\xi_1,\xi_2$. Modulo a total derivative with weight $W$, the relation is
\[
 \xi_1+\xi_2+\kappa=0.
\]
Set $\xi_1=\xi$, $\xi_2=-\xi-\kappa$, and
\begin{equation}\label{ag:eq:eta}
 \eta=-\xi(\xi+\kappa).
\end{equation}
The symbol of $(v^{(r)})^2$ is $\eta^r$, while the symbol of $vv'$ is
$-\kappa/2$. A weighted derivative is annihilated because its symbol is
multiplied by $\xi_1+\xi_2+\kappa$.

\begin{lemma}[Uniqueness of the square symbol]\label{ag:lem:symbol}
If a finite sum $W\sum_{r=0}^M c_r(v^{(r)})^2$ is a weighted total
derivative of a constant-coefficient quadratic differential expression,
then every $c_r$ is zero. The same conclusion holds coefficientwise for
formal series in $t$.
\end{lemma}
\begin{proof}
Its symbol is $\sum c_r\eta^r=0$. Since $\eta$ is a polynomial of
degree two in $\xi$, distinct powers have different degrees and are
linearly independent. The formal statement follows at every finite
coefficient of $t$. Conversely, the symmetric symbols invariant under
$\xi\mapsto-\xi-\kappa$ form the polynomial ring in $\eta$, so this
symbol also supplies the usual diagonal reduction of quadratic forms.
\end{proof}

Lemma~\ref{ag:lem:appell-primitive} ensures that this is exactly the
quotient needed after \eqref{ag:eq:affine}: the quadratic part of the
primitive of $\mathcal C-\NF\mathcal C$ is $W$ times an expression
whose coefficients do not depend on $x$. An arbitrary $x$-dependent
antiderivative, which would invalidate the symbol argument, is not
being allowed.

\subsection{The symbol of the centered Gamma exponential}
Define the even differential series
\[
 K_t(z)=\frac{2\sinh(tz/2)}{z},\qquad K_t(0)=t.
\]
The exponent in \eqref{ag:eq:centered} is $K_t(\D)u_0$. Its affine
part is $t(q+px)$, so the quadratic perturbation of $\mathcal C$ is
\begin{equation}\label{ag:eq:quadratic-density}
 [\varepsilon^2]\mathcal C=\frac{W}{2}\bigl(K_t(\D)v\bigr)^2.
\end{equation}
Now set $y=t^2p/2$ and $h=t(\xi+\kappa/2)$; then
$h^2=y^2-t^2\eta$. The product-to-sum identity for hyperbolic sines
gives the exact symbol
\begin{align}
 \frac12K_t(-\xi-\kappa)K_t(\xi)
 &=\frac{\cosh y-\cosh\sqrt{y^2-t^2\eta}}{\eta}\notag\\
 &=\sum_{r\ge0}\frac{(-1)^rt^{2r+2}}{(r+1)!}
           \piy^{\,r+1}\cosh y\,\eta^r.
 \label{ag:eq:master-symbol}
\end{align}
The square root introduces no branch: $\cosh\sqrt z$ means its entire
series $\sum_{n\ge0}z^n/(2n)!$. The numerator is divisible by $\eta$.

On the other hand, normal support and the Appell factor imply
\eqref{ag:eq:normal-sector}, initially with unknown $f_0,g_r$.
No normal monomial has exactly one factor of index at least two.
If it has exactly two, their indices must coincide. All other such
monomials contribute only at order $\varepsilon^3$ or higher.
The quadratic perturbation of its low-index part is
\[
 W\left\{\frac{t^2f_0}{2}v^2+\frac{t^3f_0'}{2}vv'
                         +\frac{t^4f_0''}{8}(v')^2\right\}.
\]
After replacing the mixed term by its symbol, the coefficients of
$\eta^0$ and $\eta^1$ are
\[
 \frac{t^2}{2}(f_0-yf_0'),\qquad \frac{t^4}{8}f_0''.
\]
For $r\ge2$ the coefficient of $\eta^r$ is simply
$t^{2r+2}g_r(y)$. Lemma~\ref{ag:lem:symbol} and
\eqref{ag:eq:master-symbol} therefore prove the derivative formula
for $g_r$ in Theorem~\ref{ag:thm:main}.

The coefficient of $\eta^0$ gives
\begin{equation}\label{ag:eq:f0-ode}
 f_0-yf_0'=\frac{\sinh y}{y}.
\end{equation}
The normal form excludes a term linear in $u_1$, hence a term linear
in $y$. With $f_0(0)=1$, the unique solution is
\eqref{ag:eq:f0-intro}. The $\eta^1$ equation agrees with its second
derivative and supplies an independent normalization check.

\subsection{The beta representation and explicit coefficients}
\begin{proposition}[Entire kernels and their recurrences]\label{ag:prop:kernel}
For every $r\ge0$ define
$g_r(y)=(-1)^r\piy^{r+1}\cosh y/(r+1)!$. Then
\begin{equation}\label{ag:eq:beta-kernel}
 g_r(y)=c_r\int_0^1(1-v^2)^r\cosh(yv)\dd v,
 \qquad c_r=\frac{(-1)^r}{2^{2r+1}r!(r+1)!}.
\end{equation}
They are even entire functions and satisfy
\begin{align}
 \sum_{r\ge0}g_r(y)w^r
   &=\frac{\cosh y-\cosh\sqrt{y^2-w}}{w},
       \label{ag:eq:bivariate}\\
 g_{r+1}(y)&=-\frac1{r+2}\piy g_r(y),
       \label{ag:eq:kernel-rec}\\
 yg_r''+(2r+2)g_r'-yg_r&=0,
 \qquad g_r(0)=\frac{(-1)^r}{(2r+2)!}.
       \label{ag:eq:kernel-ode}
\end{align}
In \eqref{ag:eq:bivariate}, the terms $r=0,1$ are symbol coefficients;
they are not $f_0$ and a nonexistent normal $u_1$ term.
\end{proposition}
\begin{proof}
Expanding the hyperbolic cosine in \eqref{ag:eq:beta-kernel} is uniformly
convergent on compact sets. For every even $k\ge0$,
\[
 \int_0^1v^k(1-v^2)^r\dd v
 =\frac12 B\left(\frac{k+1}{2},r+1\right)
 =\frac{2^r r!}{\prod_{j=0}^r(k+2j+1)}.
\]
Thus
\begin{equation}\label{ag:eq:kernel-coeff}
 k![y^k]g_r(y)=
 \frac{(-1)^r}{2^{r+1}(r+1)!\prod_{j=0}^r(k+2j+1)}
 \quad(k\text{ even}),
\end{equation}
and odd coefficients vanish. Comparing with the differentiated cosine
series proves the beta identity and the ${}_0F_1$ expression.
Taylor expansion in $w$ proves \eqref{ag:eq:bivariate}; differentiation
proves \eqref{ag:eq:kernel-rec}. Substituting the coefficients proves
\eqref{ag:eq:kernel-ode}. Finally, expand $e^{tu_0}$ and set
$y=u_1t^2/2$ to obtain \eqref{ag:eq:coefficient}.
\end{proof}

The first two high-index kernels are elementary:
\begin{align}
 g_2(y)&=\frac{(y^2+3)\sinh y-3y\cosh y}{48y^5},
       \label{ag:eq:g2}\\
 g_3(y)&=-\frac{(y^3+15y)\cosh y-(6y^2+15)\sinh y}{384y^7}.
       \label{ag:eq:g3}
\end{align}
At the origin these formulas should be evaluated by their even Taylor
series or by ${}_0F_1$, rather than by subtracting nearly equal floating
numbers. The recurrence generates every higher elementary hyperbolic
expression, while the beta form is preferable for the analytic bridge.

% END sections/03_polarization.tex


% BEGIN sections/04_profiles.tex
\section{Symmetric profiles and mixed Gamma ratios}
\label{ag:sec:profiles}

The affine argument does not depend on the particular uniform interval
average in the Gamma ratio. Let $K(z)\in\C[[z^2]]$ be even, with
$K(0)=1$, and form
\begin{equation}\label{ag:eq:profile-density}
 \mathcal C_K(t)=\exp\bigl(tK(t\D)u_0\bigr).
\end{equation}
Its coefficients are differential polynomials, and its normal form again
has an Appell factor. Denote the coefficient of $u_r^2t^{2r+2}$ after
removing that factor by $g_{K,r}(y)$, with $y=u_1t^2/2$.

\begin{theorem}[The symmetric-profile coefficient theorem]
\label{ag:thm:profile}
For every $r\ge2$,
\begin{equation}\label{ag:eq:profile-symbol}
 \boxed{
 g_{K,r}(y)=[w^r]\frac12
 K\left(\sqrt{y^2-w}-y\right)
 K\left(\sqrt{y^2-w}+y\right).}
\end{equation}
The product is even in the square-root variable, so its interpretation
as a formal series in $y,w$ is unambiguous. Formula
\eqref{ag:eq:profile-symbol} determines all coefficients of
$u_0^a u_1^k u_r^2$ in $\NF\mathcal C_K$.
\end{theorem}
\begin{proof}
Under \eqref{ag:eq:affine}, the quadratic term is
$t^2W(K(t\D)v)^2/2$. Its symbol is
$t^2K(h-y)K(h+y)/2$, where $h^2=y^2-t^2\eta$.
The primitive has Appell covariance by
Lemma~\ref{ag:lem:appell-primitive}, and normal support excludes any
other contribution to $\eta^r$ for $r\ge2$. Comparing coefficients
of $t^{2r+2}\eta^r$ proves the formula. Evenness of $K$ makes the
product invariant under $h\mapsto-h$, justifying its formal meaning.
\end{proof}

\subsection{A finite collection of Gamma scales}
For nonzero scales $\rho_1,\ldots,\rho_d$ and constants $c_i$ with
$\sum c_i=1$, put
\[
 K(z)=\sum_{i=1}^dc_iK_{\rho_i}(z),\qquad
 K_\rho(z)=\frac{2\sinh(\rho z/2)}{\rho z}.
\]
Under the digamma substitution, \eqref{ag:eq:profile-density} becomes
\begin{equation}\label{ag:eq:mixed-Gamma}
 \prod_{i=1}^d
 \left(\frac{\Gamma(x+\rho_it/2)}{\Gamma(x-\rho_it/2)}\right)^{c_i/\rho_i}.
\end{equation}
The powers are defined by the analytic logarithm which vanishes at
$t=0$. For positive real scales, real $x>0$, and sufficiently small
real $t$, this is the usual positive real power.

For compactness write
$D_n=\partial_{y^2}^{\,n}$,
$A_{ij}(y)=\cosh((\rho_i-\rho_j)y/2)$, and
$B_{ij}(y)=\cosh((\rho_i+\rho_j)y/2)$.
Theorem~\ref{ag:thm:profile} gives the explicit mixed-scale formula
\begin{align}
 g_{K,r}(y)
 =\frac{(-1)^r}{(r+1)!}\bigg\{&
 \sum_i\frac{c_i^2}{\rho_i^2}D_{r+1}\cosh(\rho_i y)\notag\\
 &+2\sum_{i<j}\frac{c_ic_j}{\rho_i\rho_j}
 \bigl(A_{ij}D_{r+1}B_{ij}-B_{ij}D_{r+1}A_{ij}\bigr)\bigg\}.
 \label{ag:eq:mixed-square}
\end{align}
Indeed, expand the product in \eqref{ag:eq:profile-symbol} and use the
hyperbolic product-to-sum identities before Taylor expanding in $w$.
Coincident scales have a regular limit; no division by
$\rho_i-\rho_j$ appears.

As one concrete check, take $c_1=c_2=1/2$, $\rho_1=1$, and $\rho_2=2$.
The coefficient of $u_2^2$ in the weight-six exponential coefficient,
normalized by $6!$, is
\begin{equation}\label{ag:eq:mixed-example}
 [u_2^2]\,6![t^6]\NF\mathcal C_K(t)=\frac{227}{32}.
\end{equation}
It can also be verified without the symbol: the exponent is
$tu_0+(5/48)t^3u_2+(17/3840)t^5u_4+\cdots$, and
$\NF(u_0u_4)=u_2^2$. The two contributions are $51/16$ and $125/32$.

The general theorem is useful even when $K$ is a polynomial rather
than a Gamma profile. Its proof does not use analytic convergence of
$K$. For convergent profiles the resulting kernels may have additional
analytic representations; the remainder of this article treats the
uniform Gamma profile, whose beta weight gives particularly short
endpoint formulas.

% END sections/04_profiles.tex


% BEGIN sections/05_hurwitz.tex
\section{A finite-endpoint Mellin--Hurwitz transform}
\label{ag:sec:hurwitz}

Fix $r\ge2$ and initially take real $a>b>0$. Define
\begin{equation}\label{ag:eq:Hmellin}
 \HH_r(s;a,b)=\frac1{\Gamma(s)}\int_0^\infty
       \frac{t^{s-1}e^{-at}g_r(bt)}{1-e^{-t}}\dd t,
                  \qquad \Re s>1.
\end{equation}
At zero the integrand is $O(t^{\Re s-2})$, and at infinity the
exponential decay is governed by $a-b>0$. These conditions justify
ordinary integration and all fixed-order spectral derivatives in the
stated half-plane.

The Hurwitz zeta convention is
$\zeta(s,x)=\sum_{n\ge0}(n+x)^{-s}$ for $\Re s>1$, continued
meromorphically in $s$. Its standard Mellin representation and spatial
derivative are recorded in \cite{DLMFzeta}. They also follow directly
from the series in their common domain.

\begin{theorem}[The beta average and its finite endpoints]
\label{ag:thm:hurwitz}
With $c_r$ as in \eqref{ag:eq:beta-kernel},
\begin{equation}\label{ag:eq:Hcompact}
 \HH_r(s;a,b)=\frac{c_r}{2}\int_{-1}^1
                    (1-v^2)^r\zeta(s,a+bv)\dd v.
\end{equation}
This provides meromorphic continuation to every $s$, with the single
possible pole $s=1$, whose residue is
\begin{equation}\label{ag:eq:Hresidue}
 g_r(0)=\frac{(-1)^r}{(2r+2)!}.
\end{equation}
Define
\begin{equation}\label{ag:eq:weights}
 w_{rj}(b)=\frac{(-1)^j j!\binom r{j-r}}
 {2^{j+2}r!(r+1)!b^{j+1}},\qquad r\le j\le2r.
\end{equation}
Then, as an identity of meromorphic functions,
\begin{equation}\label{ag:eq:Hendpoints}
 \boxed{
 \HH_r(s;a,b)=\sum_{j=r}^{2r}w_{rj}(b)
 \frac{\zeta(s-j-1,a+b)-(-1)^j\zeta(s-j-1,a-b)}{(1-s)_{j+1}}.}
\end{equation}
At apparent poles of the separate summands, the entire sum is to be
continued jointly. No pole other than $s=1$ survives.
\end{theorem}
\begin{proof}
Insert \eqref{ag:eq:beta-kernel} into \eqref{ag:eq:Hmellin}.
Absolute convergence permits interchanging the integrals. Combining
the two exponentials in the hyperbolic cosine yields
\eqref{ag:eq:Hcompact}. On every compact subset of $\Re x>0$,
$\zeta(s,x)-1/(s-1)$ is holomorphic and locally bounded in $s$.
Its compact $v$-integral is therefore entire, proving the continuation
and the residue. The integral of the weight is the one already used
in \eqref{ag:eq:kernel-coeff}.

For the endpoint identity let $P(v)=(1-v^2)^r$. Repeated spatial
antidifferentiation gives
\[
 \frac{\mathrm d^{j+1}}{\mathrm dv^{j+1}}
 \frac{\zeta(s-j-1,a+bv)}{b^{j+1}(1-s)_{j+1}}
 =\zeta(s,a+bv),
\]
initially for $s$ outside the finite exceptional set. Integrating by
parts until the polynomial vanishes leaves
\[
 \sum_{j=0}^{2r}\frac{(-1)^j}{b^{j+1}(1-s)_{j+1}}
       [P^{(j)}(v)\zeta(s-j-1,a+bv)]_{v=-1}^{v=1}.
\]
The derivatives with $j<r$ vanish at both endpoints, and for $r\le j\le2r$,
\begin{equation}\label{ag:eq:weight-derivative}
 P^{(j)}(1)=(-1)^r j!\binom r{j-r}2^{2r-j},\qquad
 P^{(j)}(-1)=(-1)^jP^{(j)}(1).
\end{equation}
Multiplication by $c_r/2$ proves \eqref{ag:eq:Hendpoints}. Analytic
continuation establishes it at every remaining regular point.
\end{proof}

The formulas extend holomorphically in $a,b$ when $\Re a>|b|$, using
the right-half-plane Hurwitz branch. They are even in $b$. At $b=0$
the compact formula gives the confluent value
\begin{equation}\label{ag:eq:Hcollision}
 \HH_r(s;a,0)=g_r(0)\zeta(s,a).
\end{equation}
The negative powers of $b$ in \eqref{ag:eq:Hendpoints} must cancel;
one should not substitute $b=0$ separately in those terms.

\subsection{All logarithmically weighted Mellin integrals}
Let $\mathsf Y_m$ be the complete exponential Bell polynomial, defined
by $\exp(\sum_{j\ge1}x_jz^j/j!)=\sum_{m\ge0}\mathsf Y_m(x_1,\ldots,x_m)z^m/m!$.
Differentiating the normalized Mellin factor yields, for $\Re s>1$,
\begin{align}
 \partial_s^m\HH_r(s;a,b)
 =\frac1{\Gamma(s)}\int_0^\infty&
 \frac{t^{s-1}e^{-at}g_r(bt)}{1-e^{-t}}\notag\\[-2pt]
 &\quad\times\mathsf Y_m\bigl(\log t-\psi(s),-\psi'(s),
                      \ldots,-\psi^{(m-1)}(s)\bigr)\dd t.
 \label{ag:eq:log-Mellin}
\end{align}
For $m=0$ the Bell factor is one. For $m=1,2$ the factors are
$\log t-\psi(s)$ and $(\log t-\psi(s))^2-\psi'(s)$.
This is an exact normalization, not the unnormalized rule of merely
inserting $(\log t)^m$. Endpoint differentiation of
\eqref{ag:eq:Hendpoints} therefore gives finite evaluations of all
these convergent logarithmic integrals.

% END sections/05_hurwitz.tex


% BEGIN sections/06_stieltjes.tex
\section{Every Stieltjes order and anchored parameter primitives}
\label{ag:sec:stieltjes}

We use the Laurent convention
\begin{equation}\label{ag:eq:Stieltjes-convention}
 \zeta(1+\varepsilon,x)=\frac1\varepsilon+
       \sum_{m\ge0}\frac{(-1)^m}{m!}\gamma_m(x)\varepsilon^m,
                 \qquad \gamma_0(x)=-\psi(x).
\end{equation}
Define the ordinary weighted averages
\begin{equation}\label{ag:eq:Bdefinition}
 \BB_{r,m}(a,b)=\frac{c_r}{2}\int_{-1}^1
                      (1-v^2)^r\gamma_m(a+bv)\dd v.
\end{equation}
Consequently
\[
 \HH_r(1+\varepsilon;a,b)=\frac{g_r(0)}\varepsilon+
              \sum_{m\ge0}\frac{(-1)^m}{m!}\BB_{r,m}(a,b)\varepsilon^m.
\]
This expansion is locally uniform for the spatial parameters under the
conditions of Theorem~\ref{ag:thm:hurwitz}.

\subsection{Finite complete-harmonic coefficients}
For $j\ge1$, let
\begin{equation}\label{ag:eq:homogeneous}
 E_j(\varepsilon)=\prod_{n=1}^j(1-\varepsilon/n)^{-1}
               =\sum_{\ell\ge0}h_\ell^{(j)}\varepsilon^\ell.
\end{equation}
These are complete homogeneous polynomials in $1,1/2,\ldots,1/j$.
They can be calculated from generalized harmonic numbers without
expanding a product:
\begin{equation}\label{ag:eq:hrecurrence}
 h_0^{(j)}=1,\qquad
 \ell h_\ell^{(j)}=\sum_{d=1}^\ell H_j^{(d)}h_{\ell-d}^{(j)}.
\end{equation}
For example, $h_1^{(j)}=H_j$ and
$h_2^{(j)}=((H_j)^2+H_j^{(2)})/2$.

\begin{theorem}[All-order Stieltjes endpoint formula]
\label{ag:thm:Stieltjes}
Let $D_{rj}=w_{rj}(b)/j!$, and use $\zeta^{(k)}$ for a derivative
in the spectral variable. For every $m\ge0$,
\begin{align}
 \boxed{\displaystyle
 \BB_{r,m}(a,b)=(-1)^{m+1}m!
 \sum_{j=r}^{2r}D_{rj}\sum_{k=0}^{m+1}
 \frac{h_{m+1-k}^{(j)}}{k!}
 \left\{\zeta^{(k)}(-j,a+b)-(-1)^j\zeta^{(k)}(-j,a-b)\right\}.}
 \label{ag:eq:Stieltjes-all}
\end{align}
There are only $r+1$ endpoint pairs, and the largest spectral derivative
order is $m+1$.
\end{theorem}
\begin{proof}
At $s=1+\varepsilon$,
\[
 (1-s)_{j+1}=-j!\varepsilon\prod_{n=1}^j(1-\varepsilon/n).
\]
Set
$F_j(\varepsilon)=\zeta(-j+\varepsilon,a+b)
 -(-1)^j\zeta(-j+\varepsilon,a-b)$.
Theorem~\ref{ag:thm:hurwitz} becomes
\begin{equation}\label{ag:eq:H-Laurent-master}
 \HH_r(1+\varepsilon;a,b)
 =-\frac1\varepsilon\sum_{j=r}^{2r}D_{rj}
                          E_j(\varepsilon)F_j(\varepsilon).
\end{equation}
The constant coefficient of the numerator is $-g_r(0)$.
The coefficient of $\varepsilon^{m+1}$ is a finite convolution of
\eqref{ag:eq:homogeneous} with the Taylor series of $F_j$.
Comparison with \eqref{ag:eq:Stieltjes-convention} proves
\eqref{ag:eq:Stieltjes-all}. Taking a logarithmic derivative of
\eqref{ag:eq:homogeneous} proves \eqref{ag:eq:hrecurrence}.
\end{proof}

At $m=0$ the inner expression is $F_j'(0)+H_jF_j(0)$.
At $m=1$ it is
$F_j''(0)/2+H_jF_j'(0)+((H_j)^2+H_j^{(2)})F_j(0)/2$.
Thus the harmonic counterterms are required even though every spatial
integral in \eqref{ag:eq:Bdefinition} is nonsingular. They arise from
the finite endpoint denominators, not from an arbitrary regularization.

\subsection{Spatial derivatives}
Differentiation of the compact formula gives, for $L\ge0$,
\begin{equation}\label{ag:eq:H-spatial}
 \partial_a^L\HH_r(s;a,b)=(-1)^L(s)_L\HH_r(s+L;a,b).
\end{equation}
For $L\ge1$, comparison at $s=1+\varepsilon$ yields
\begin{equation}\label{ag:eq:B-spatial}
 \partial_a^L\BB_{r,m}(a,b)
 =(-1)^{m+L}m![\varepsilon^m]
       (1+\varepsilon)_L\HH_r(1+L+\varepsilon;a,b).
\end{equation}
In particular,
\[
 \partial_a\BB_{r,m}=(-1)^{m+1}
  \left\{\partial_s^m\HH_r(2;a,b)
                   +m\partial_s^{m-1}\HH_r(2;a,b)\right\},
\]
where the second term is omitted for $m=0$. The finite product
$(1+\varepsilon)_L$ gives another explicit generalized-harmonic
coefficient expansion when desired.

\subsection{First primitives and the Gamma endpoint}
\begin{corollary}[An anchored first primitive]\label{ag:cor:primitive}
For $a,a_0>b>0$ and $m\ge0$,
\begin{equation}\label{ag:eq:first-primitive}
 \int_{a_0}^a\BB_{r,m}(\alpha,b)\dd\alpha
 =\frac{(-1)^{m+1}}{m+1}
  \left.\partial_s^{m+1}\bigl(\HH_r(s;a,b)-\HH_r(s;a_0,b)\bigr)\right|_{s=0}.
\end{equation}
For $m=0$ this is the ordinary log-Gamma identity
\begin{equation}\label{ag:eq:logGamma-primitive}
 \int_{a_0}^a\BB_{r,0}(\alpha,b)\dd\alpha
 =-\frac{c_r}{2}\int_{-1}^1(1-v^2)^r
       \log\frac{\Gamma(a+bv)}{\Gamma(a_0+bv)}\dd v.
\end{equation}
\end{corollary}
\begin{proof}
Equation \eqref{ag:eq:H-spatial} with $L=1$ gives
$\partial_a\HH_r(s)=-s\HH_r(1+s)$. Its coefficient of $s^{m+1}$ is
$(-1)^{m+1}\BB_{r,m}/m!$, proving
\eqref{ag:eq:first-primitive} after differentiation and anchoring.
Use $\zeta'(0,x)=\log\Gamma(x)-\tfrac12\log(2\pi)$
\cite{DLMFzeta} in \eqref{ag:eq:Hcompact}; the constant cancels between
the two endpoints, giving \eqref{ag:eq:logGamma-primitive}.
\end{proof}

\subsection{An explicit primitive of every depth}
The same argument can be made completely normalized at arbitrary
integration depth, without introducing unspecified constants.
For $L\ge1$ set
\[
 U_L(s;a,b)=\frac{\HH_r(s-L;a,b)}{(1-s)_L}.
\]
Define
\begin{align}
 V_L(s;a,a_0,b)={}&U_L(s;a,b)
 -\sum_{j=0}^{L-1}\frac{(a-a_0)^j}{j!}\partial_a^jU_L(s;a_0,b)\notag\\
 &-\frac{g_r(0)(a-a_0)^L}{L!(s-1)}.
 \label{ag:eq:anchored-L}
\end{align}
\begin{theorem}[All anchored Stieltjes primitives]\label{ag:thm:all-primitives}
The function $V_L$ is holomorphic near $s=1$, and
\begin{equation}\label{ag:eq:all-primitives}
 J_{L,m}(a;a_0,b)=(-1)^m m![\varepsilon^m]
                          V_L(1+\varepsilon;a,a_0,b)
\end{equation}
satisfies
$\partial_a^LJ_{L,m}=\BB_{r,m}(a,b)$ and
$\partial_a^jJ_{L,m}(a_0;a_0,b)=0$ for $0\le j<L$.
It is the unique $L$-fold primitive with those anchor conditions.
\end{theorem}
\begin{proof}
The identity $\partial_a^LU_L=\HH_r(s;a,b)$ follows from
\eqref{ag:eq:H-spatial}. The residue of $U_L$ at $s=1$ is a polynomial
of degree at most $L$ in $a$, because
$\zeta(1-L,a+bv)=-B_L(a+bv)/L$. Its $L$th derivative is $g_r(0)$.
Taylor subtraction through degree $L-1$ therefore leaves exactly the
last pole written in \eqref{ag:eq:anchored-L}; removing it proves
holomorphy. Differentiating $L$ times leaves
$\HH_r(s;a,b)-g_r(0)/(s-1)$. Coefficient comparison proves the
claimed derivative. All lower anchor derivatives vanish by construction.
The difference of any two such primitives is a polynomial of degree
below $L$ with all its anchor derivatives zero, hence vanishes.
\end{proof}

% END sections/06_stieltjes.tex


% BEGIN sections/07_polylogs.tex
\section{Lerch deformation and finite polylogarithmic closure}
\label{ag:sec:polylogs}

For $|z|<1$, replace the denominator in \eqref{ag:eq:Hmellin} by
$1-ze^{-t}$ and define
\begin{equation}\label{ag:eq:Lerch-transform}
 \HH_r(s;a,b;z)=\frac1{\Gamma(s)}\int_0^\infty
             \frac{t^{s-1}e^{-at}g_r(bt)}{1-ze^{-t}}\dd t,
                     \qquad\Re s>0.
\end{equation}
Let $\Phi(z,s,x)=\sum_{n\ge0}z^n(n+x)^{-s}$ in the right half-plane.
The Lerch series, its Mellin formula, and its relation to polylogarithms
are classical \cite{DLMFlerch}. Termwise spatial differentiation gives
$\partial_x\Phi(z,s,x)=-s\Phi(z,s+1,x)$.

\begin{theorem}[Lerch and polylogarithmic endpoints]\label{ag:thm:Lerch}
For $|z|<1$,
\begin{align}
 \HH_r(s;a,b;z)&=\frac{c_r}{2}\int_{-1}^1
                       (1-v^2)^r\Phi(z,s,a+bv)\dd v,
                       \label{ag:eq:Lerch-compact}\\
 &=\sum_{j=r}^{2r}w_{rj}(b)
 \frac{\Phi(z,s-j-1,a+b)-(-1)^j\Phi(z,s-j-1,a-b)}{(1-s)_{j+1}}.
                       \label{ag:eq:Lerch-endpoint}
\end{align}
The function is entire in $s$. Hence every apparent integer pole on
the right is removable, including $s=1$.
If $a\pm b$ are positive integers, every endpoint is a finite
polylogarithmic expression, using
\begin{equation}\label{ag:eq:integer-Lerch}
 \Phi(z,s,N)=z^{-N}\left(\Li_s(z)-\sum_{n=1}^{N-1}\frac{z^n}{n^s}\right)
 \quad(N\ge1,
 z\ne0).
\end{equation}
All values at $z=0$ are supplied by analytic limits.
\end{theorem}
\begin{proof}
The beta representation and absolute convergence prove
\eqref{ag:eq:Lerch-compact}. The endpoint argument in
Theorem~\ref{ag:thm:hurwitz} applies unchanged. Exponential decay in
$n$ makes the Lerch series locally uniformly convergent for every
$s$, including all spectral derivatives. This proves entire continuation
and removability of the displayed poles. Shifting the summation index
proves \eqref{ag:eq:integer-Lerch}.
\end{proof}

The limit $z\to1$ is initially justified for $\Re s>1$, where it
recovers the Hurwitz transform. Continuing in $s$ before taking that
limit need not be interchangeable at $s=1$: the Hurwitz pole must be
kept. At a removable integer in \eqref{ag:eq:Lerch-endpoint}, the value
may involve \emph{order derivatives} of polylogarithms; a formula
containing a vanishing denominator should not be read as an ordinary
substitution rule.

\subsection{Rational spatial endpoints}
For $0<z<1$, let $w=z^{1/q}$, $\omega=e^{2\pi\ii/q}$, and take
positive integers $A,q$. The roots-of-unity filter in the absolutely
convergent polylogarithm series gives
\begin{align}
 \Phi(z,s,A/q)
 ={}&q^{s-1}w^{-A}\sum_{\ell=0}^{q-1}\omega^{-\ell A}
                                      \Li_s(w\omega^\ell)\notag\\
 &-q^sw^{-A}
       \sum_{\substack{1\le n<A\\n\equiv A\;({\rm mod}\ q)}}
                                    \frac{w^n}{n^s}.
 \label{ag:eq:rational-Lerch}
\end{align}
Indeed the finite Fourier sum retains exactly the positive integers
congruent to $A$ modulo $q$, and the finite prefix removes those below
$A$. The formula holds for every complex $s$ by normal convergence.
Its order derivatives give finite log-weighted polylogarithm reductions
for all spectral derivatives of the rational-endpoint transform.

\subsection{An explicit expression with only three polylogarithms}
Put
\[
 M(t)=(t^2+3)\sinh t-3t\cosh t,
\]
and, for $z\ne0$,
\[
 J_\sigma^\pm(z)=(z^{-3}\pm z^{-1})\Li_\sigma(z)
                                  -z^{-2}-2^{-\sigma}z^{-1}.
\]
For $r=2,a=2,b=1$, \eqref{ag:eq:Lerch-endpoint} reads
\begin{equation}\label{ag:eq:r2-Lerch}
 \HH_2(s;2,1;z)=\frac{J_{s-3}^-(z)}{96(1-s)_3}
       -\frac{J_{s-4}^+(z)}{32(1-s)_4}
       +\frac{J_{s-5}^-(z)}{32(1-s)_5}.
\end{equation}
At $s=6$ there is no exceptional denominator. Since
$g_2(t)=M(t)/(48t^5)$, this gives the ordinary convergent identity
\begin{align}
 \int_0^\infty\frac{e^{-2t}M(t)}{1-ze^{-t}}\dd t
 ={}&\frac4{z^2}+\frac5{4z}
 -\frac{1-z^2}{z^3}\left(\Li_3(z)+\frac32\Li_1(z)\right)\notag\\
 &-\frac{3(1+z^2)}{2z^3}\Li_2(z),\qquad 0<z<1.
 \label{ag:eq:three-polylogs}
\end{align}
The finite limit at $z=1$ is $21/4-\pi^2/2$. The individual negative
powers of $z$ cancel at $z=0$. These endpoint cancellations are most
transparent in the integral or the Lerch series, rather than in the
expanded polylogarithmic expression.

% END sections/07_polylogs.tex


% BEGIN sections/08_examples.tex
\section{Quartic weights: polygamma, Stieltjes, and hyperbolic integrals}
\label{ag:sec:examples}

Take $r=2,a=2,b=1$ and define $W(x)=(x-1)^2(3-x)^2$. Then
\begin{equation}\label{ag:eq:r2-compact}
 \HH_2(s;2,1)=\frac1{768}\int_1^3W(x)\zeta(s,x)\dd x.
\end{equation}
Integer translation of the Hurwitz zeta in
\eqref{ag:eq:Hendpoints} gives
\begin{align}
 \HH_2(s;2,1)={}&-\frac{1+2^{3-s}}{96(1-s)_3}
 -\frac{2\zeta(s-4)-1-2^{4-s}}{32(1-s)_4}
 -\frac{1+2^{5-s}}{32(1-s)_5}.
 \label{ag:eq:r2-Riemann}
\end{align}
This is a three-term meromorphic identity; the right side is not to be
evaluated term by term at its apparent poles.

\begin{proposition}[Six ordinary polygamma integrals]
\label{ag:prop:six-integrals}
The following identities hold:
\begin{align}
 \int_1^3W(x)\psi(x)\dd x
   &=\frac{83}{180}-8\zeta'(-3),\label{ag:eq:Ipsi0}\\
 \int_1^3W(x)\psi'(x)\dd x
   &=\frac{6\zeta(3)}{\pi^2},\label{ag:eq:Ipsi1}\\
 \int_1^3W(x)\psi''(x)\dd x
   &=-48\zeta'(-1)+8\log2-14,\label{ag:eq:Ipsi2}\\
 \int_1^3W(x)\psi^{(3)}(x)\dd x
   &=24\log(4\pi)-60,\label{ag:eq:Ipsi3}\\
 \int_1^3W(x)\psi^{(4)}(x)\dd x
   &=48\gamma+24\log2-46,\label{ag:eq:Ipsi4}\\
 \int_1^3W(x)\psi^{(5)}(x)\dd x
   &=84-8\pi^2.\label{ag:eq:Ipsi5}
\end{align}
\end{proposition}
\begin{proof}
Take the Laurent coefficient at $s=1$ and the regular values at
$s=2,3,4,5,6$ in \eqref{ag:eq:r2-Riemann}, respectively. Use
$\gamma_0=-\psi$ and
$\psi^{(m)}(x)=(-1)^{m+1}m!\zeta(m+1,x)$ for $m\ge1$
\cite{DLMFpsi}. To make the cancellation reproducible, the required
intermediate values are
\begin{align*}
 \HH_2(1+\varepsilon;2,1)
 &=\frac1{720\varepsilon}+\frac{\zeta'(-3)}{96}
                               -\frac{83}{138240}+O(\varepsilon),\\
 768\HH_2(2;2,1)&=-24\zeta'(-2),\\
 768\HH_2(3;2,1)&=24\zeta'(-1)-4\log2+7,\\
 768\HH_2(4;2,1)&=-8\zeta'(0)-10+4\log2,\\
 768\HH_2(5;2,1)&=-2\gamma-\log2+\frac{23}{12},\\
 768\HH_2(6;2,1)&=\frac7{10}-\frac{\pi^2}{15}.
\end{align*}
Finally, $\zeta'(0)=-\tfrac12\log(2\pi)$ and
$\zeta'(-2)=-\zeta(3)/(4\pi^2)$ give the printed forms.
The latter is the derivative of the zeta functional equation at its
simple trivial zero; see also \cite{DLMFintegers}.
\end{proof}

The next Stieltjes coefficient is equally explicit:
\begin{equation}\label{ag:eq:gamma1-example}
 \boxed{
 \int_1^3W(x)\gamma_1(x)\dd x
 =\frac{367}{432}-\frac43\log2
                 -\frac{44}{3}\zeta'(-3)-4\zeta''(-3).}
\end{equation}
To verify it algebraically, put
$\zeta(-3+\varepsilon)=1/120+\zeta'(-3)\varepsilon+
\zeta''(-3)\varepsilon^2/2+O(\varepsilon^3)$ and
$2^{-\varepsilon}=1-\varepsilon\log2+
\varepsilon^2(\log2)^2/2+O(\varepsilon^3)$ in
\eqref{ag:eq:r2-Riemann}. The coefficient of $\varepsilon$ is minus
$1/768$ times the left side of \eqref{ag:eq:gamma1-example}; the
logarithm-square terms cancel. Theorem~\ref{ag:thm:Stieltjes} proves the
corresponding finite identity at every higher order.

\subsection{A convergent hyperbolic integral}
At $s=2$, \eqref{ag:eq:Hmellin} and \eqref{ag:eq:g2} also give
\begin{equation}\label{ag:eq:hyperbolic-integral}
 \boxed{
 \int_0^\infty
 \frac{e^{-2t}\bigl((t^2+3)\sinh t-3t\cosh t\bigr)}
      {t^4(1-e^{-t})}\dd t
 =\frac{3\zeta(3)}{8\pi^2}.}
\end{equation}
The numerator in parentheses is $t^5/15+O(t^7)$ at zero, so the
integrand has a finite limit. At infinity it decays exponentially
with a rational power of $t$. There is no hidden finite part.
Multiplying the value of $\HH_2(2;2,1)$ by $48$ proves the identity.

These examples show two uses of the same kernel: its normal-form
coefficients determine a Gamma differential sector, while its beta
representation evaluates an independent family of ordinary integrals.
They should not be conflated with a claim that all mixed Gamma moments
have thereby been evaluated.

% END sections/08_examples.tex


% BEGIN sections/09_bernoulli.tex
\section{Bernoulli selectors: single-zeta closure and a classification}
\label{ag:sec:bernoulli}

The beta weight in \eqref{ag:eq:Hcompact} is even and vanishes with its
first derivative at both endpoints. A different basis of precisely this
polynomial space makes the arithmetic of the interval $[1,3]$ diagonal.
Define, for $r\ge2$,
\begin{equation}\label{ag:eq:Bernoulli-weight}
 P_r(v)=B_{2r}\left(\frac{v+1}{2}\right)-B_{2r},\qquad
 I_r(s)=\int_1^3P_r(x-2)\zeta(s,x)\dd x.
\end{equation}
Here $B_n(x)$ is defined by
$te^{xt}/(e^t-1)=\sum_{n\ge0}B_n(x)t^n/n!$, with $B_n=B_n(0)$
\cite{DLMFBernoulli}. In particular $B_1=-1/2$ but $B_1(1)=1/2$.
The endpoint derivatives are
\begin{equation}\label{ag:eq:Bernoulli-d}
 d_{rj}:=P_r^{(j)}(1)
 =\frac{(2r)!}{2^j(2r-j)!}B_{2r-j}(1),\qquad 1\le j\le2r.
\end{equation}
The only nonzero odd derivative in this range is $j=2r-1$.

\subsection{The full complex-order selector identity}
\begin{theorem}[One shifted zeta at every complex order]
\label{ag:thm:single-zeta}
As meromorphic functions of $s$,
\begin{equation}\label{ag:eq:single-zeta}
 \boxed{
 I_r(s)=-\frac{(2r)!}{2^{2r-1}(1-s)_{2r}}\zeta(s-2r)
 -\sum_{j=2}^{2r}\frac{(-1)^jd_{rj}}{(1-s)_{j+1}}
                                    \bigl(1+2^{j+1-s}\bigr).}
\end{equation}
The compact integral has only the pole $s=1$, of residue $-2B_{2r}$;
all other apparent poles in the display cancel.
\end{theorem}
\begin{proof}
The Bernoulli reflection identity makes $P_r$ even. The endpoint values
and first derivatives vanish. Repeated integration by parts, exactly as
in Theorem~\ref{ag:thm:hurwitz}, gives
\[
 I_r(s)=\sum_{j=2}^{2r}\frac{(-1)^jd_{rj}}{(1-s)_{j+1}}
       \bigl\{\zeta(s-j-1,3)-(-1)^j\zeta(s-j-1,1)\bigr\}.
\]
Integer translation of Hurwitz zeta makes the braces
\[
 (1-(-1)^j)\zeta(s-j-1)-1-2^{j+1-s}.
\]
All even $j$ therefore contribute only elementary terms. For odd $j$,
\eqref{ag:eq:Bernoulli-d} vanishes except when $j=2r-1$, because all
odd Bernoulli numbers beyond $B_1$ are zero. At this last index,
$d_{r,2r-1}=(2r)!/2^{2r}$. This gives
\eqref{ag:eq:single-zeta}. Finally,
$\int_{-1}^1P_r(v)\dd v=-2B_{2r}$, since the integral of
$B_{2r}(x)$ over $[0,1]$ is zero. The compact representation proves the
pole assertion.
\end{proof}

\begin{theorem}[The selector basis and its uniqueness]
\label{ag:thm:selector-classification}
Let $\mathscr V_R$ be the rational vector space of even polynomials
$P(v)$ of degree at most $2R$ such that $P(1)=P'(1)=0$, where $R\ge2$.
Then $P_2,\ldots,P_R$ form a basis, with the explicit expansion
\begin{equation}\label{ag:eq:selector-expansion}
 \boxed{
 P(v)=\sum_{r=2}^R\frac{2^{2r}P^{(2r-1)}(1)}{(2r)!}\,P_r(v).}
\end{equation}
Moreover, for a nonzero $P\in\mathscr V_R$ and a fixed $r\in\{2,\ldots,R\}$,
the identity
\begin{equation}\label{ag:eq:selector-class}
 \int_1^3P(x-2)\zeta(s,x)\dd x
 =F(s)\zeta(s-2r)+G(s)+2^{-s}H(s)
\end{equation}
with rational functions $F,G,H\in\C(s)$ holds if and only if $P$ is a
nonzero scalar multiple of $P_r$. For rational $P$, that scalar is
rational. This is a statement about meromorphic functions of $s$, not
about independence of particular numerical zeta values.
\end{theorem}
\begin{proof}
The dimension of the even-polynomial space is $R+1$. Evaluation of
$P(1)$ and $P'(1)$ imposes two independent linear conditions, so
$\dim\mathscr V_R=R-1$. Each $P_r$ belongs to this space and has degree
$2r$ with nonzero leading coefficient. They are therefore a basis.
For the odd derivative of order $2j-1$ at $1$, all basis elements vanish
except $P_j$, whose value is $(2j)!/2^{2j}$. This proves
\eqref{ag:eq:selector-expansion}.

It remains to justify genuine functional uniqueness, rather than only
uniqueness in a chosen formal list. The functions
\begin{equation}\label{ag:eq:functional-independence}
 1,\quad 2^{-s},\quad \zeta(s-4),\quad\zeta(s-6),\ldots,\zeta(s-2R)
\end{equation}
are linearly independent over $\C(s)$. To see this, clear rational
denominators in a purported relation. In a right half-plane its
Dirichlet expansion has the form
\[
 \sum_{n\ge1}Q_n(s)n^{-s}=0,
\]
where each $Q_n$ is a polynomial of uniformly bounded degree in $s$.
For $n\ge3$, it is
$Q_n(s)=\sum_{r=2}^Rp_r(s)n^{2r}$ for fixed polynomials $p_r$;
only $n=1,2$ have the extra polynomial terms.

If some $Q_n$ were nonzero, choose the smallest such $n_0$ and divide
by $n_0^{-s}$ for real $s\to+\infty$. The remaining tail is bounded by
a polynomial in $s$ times an exponentially decaying function. This
follows from $|Q_n(s)|\le C(1+s)^M n^{2R}$ and the convergent series
$\sum_{n>n_0}n^{2R}(n_0/n)^s$; splitting off half the exponent bounds
it by $C'(n_0/(n_0+1))^{s/2}$ for large $s$.
A nonzero polynomial $Q_{n_0}(s)$ cannot tend to zero exponentially.
Thus every $Q_n$ is zero. For infinitely many integers $n\ge3$,
$\sum p_r(s)n^{2r}=0$ as a polynomial in $s$; regarding it as a
polynomial in $n$ gives $p_r=0$ for every $r$. The two remaining
coefficients also vanish. This proves \eqref{ag:eq:functional-independence}.

Expand $P$ by \eqref{ag:eq:selector-expansion} and apply
Theorem~\ref{ag:thm:single-zeta}. The independence just proved forces
all coefficients except the prescribed $r$ to be zero in
\eqref{ag:eq:selector-class}. The converse follows from
\eqref{ag:eq:single-zeta}.
\end{proof}

The restriction on the elementary remainder in
\eqref{ag:eq:selector-class} is essential. Allowing arbitrary
meromorphic coefficients would make such a classification vacuous.
Likewise, this functional theorem says nothing about possible
unexpected numerical relations after specializing $s$.

\subsection{Beta--Bernoulli transport}
For the particular polynomial $(1-v^2)^r$, equation
\eqref{ag:eq:weight-derivative} makes the selector expansion completely
explicit:
\begin{equation}\label{ag:eq:beta-Bernoulli}
 (1-v^2)^r=(-1)^r2^{2r}\sum_{j=2}^{r}
       \frac1j\binom r{2j-r-1}
       \left(B_{2j}\left(\frac{v+1}{2}\right)-B_{2j}\right),
\end{equation}
where a binomial coefficient outside $0\le2j-r-1\le r$ is zero.
Consequently the Gamma-derived Mellin transform on $[1,3]$ satisfies
\begin{equation}\label{ag:eq:H-selector-transport}
 \HH_r(s;2,1)=\frac1{4r!(r+1)!}
                 \sum_{j=2}^{r}\frac1j\binom r{2j-r-1}I_j(s).
\end{equation}
This links the affine Gamma square kernel directly to the single-zeta
selector basis. It also explains why only certain shifted even-order
zeta coordinates appear in its integer-endpoint reduction.

\subsection{One odd zeta in each trigamma integral}
Define the rational number
\begin{equation}\label{ag:eq:Rr}
 R_r=-\sum_{j=2}^{2r}\frac{(-1)^jd_{rj}}{(j-1)!}
                         (H_{j-1}-1)(1+2^{j-1}).
\end{equation}
\begin{corollary}[Bernoulli-weighted trigamma integrals]
\label{ag:cor:odd-zeta}
For every integer $r\ge2$,
\begin{align}
 \boxed{\displaystyle
 \int_1^3\left(B_{2r}\left(\frac{x-1}{2}\right)-B_{2r}\right)
                                  \psi'(x)\dd x
 =R_r+\frac{(-1)^r(2r)!}{2^{4r-2}\pi^{2r-2}}\zeta(2r-1).}
 \label{ag:eq:odd-zeta-integral}
\end{align}
\end{corollary}
\begin{proof}
For any even $P$ with $P(\pm1)=P'(\pm1)=0$, write
$d_j=P^{(j)}(1)$. In the finite endpoint expression put
$s=2+\varepsilon$. The reciprocal denominator is
\[
 \frac1{(1-s)_{j+1}}
 =\frac1{(j-1)!\varepsilon}
     \left(1+(H_{j-1}-1)\varepsilon+O(\varepsilon^2)\right).
\]
The endpoint numerator has constant $-1-2^{j-1}$ and linear coefficient
$2\mathbf1_{j\text{ odd}}\zeta'(1-j)+2^{j-1}\log2$.
All pole terms cancel by analyticity of the compact integral.
The coefficient of $\log2$ is
\[
 \sum_{j\ge2}\frac{(-1)^j2^{j-1}d_j}{(j-1)!}=-P'(-1)=0.
\]
The remaining value is therefore
\[
 \sum_{j\ge2}\frac{(-1)^jd_j}{(j-1)!}
 \left\{2\mathbf1_{j\text{ odd}}\zeta'(1-j)
                     -(H_{j-1}-1)(1+2^{j-1})\right\}.
\]
For $P=P_r$ only $j=2r-1$ contributes a zeta derivative, and its
coefficient is $-r(2r-1)/2^{2r-2}$. Finally use
\[
 \zeta'(-2n)=\frac{(-1)^n(2n)!}{2(2\pi)^{2n}}\zeta(2n+1)
 \quad(n\ge1),
\]
obtained by differentiating the functional equation at a trivial zero
\cite{DLMFintegers}. This proves \eqref{ag:eq:odd-zeta-integral}.
\end{proof}

The first rational offsets are
\[
 R_2=0,\quad R_3=-\frac1{40},\quad R_4=\frac3{70},\quad
 R_5=-\frac{57}{560},\quad R_6=\frac{397}{1155}.
\]
For example, the $r=3$ value is
$-1/40-45\zeta(5)/(64\pi^4)$, and the $r=4$ value is
$3/70+315\zeta(7)/(128\pi^6)$.
For $r=2$ the polynomial weight is $W(x)/16$, reproducing
\eqref{ag:eq:Ipsi1}.

\subsection{All Stieltjes orders in a single negative-odd zeta jet}
The complex-order identity proves more than the trigamma specialization.
With $E_j$ from \eqref{ag:eq:homogeneous}, define
\begin{align}
 L_r(\varepsilon)={}&
 \frac{r}{2^{2r-2}}E_{2r-1}(\varepsilon)
                              \zeta(1-2r+\varepsilon)\notag\\
 &+\sum_{j=2}^{2r}\frac{(-1)^jd_{rj}}{j!}
                  E_j(\varepsilon)\bigl(1+2^je^{-\varepsilon\log2}\bigr).
 \label{ag:eq:selector-L}
\end{align}
Equation \eqref{ag:eq:single-zeta} is precisely
$I_r(1+\varepsilon)=L_r(\varepsilon)/\varepsilon$.
In particular $L_r(0)=-2B_{2r}$, and for every $m\ge0$,
\begin{equation}\label{ag:eq:selector-Stieltjes}
 \boxed{
 \int_1^3P_r(x-2)\gamma_m(x)\dd x
       =(-1)^m m![\varepsilon^{m+1}]L_r(\varepsilon).}
\end{equation}
This is a finite formula in the spectral derivatives
$\zeta^{(k)}(1-2r)$, $0\le k\le m+1$, rational numbers,
generalized harmonic numbers, and powers of $\log2$. It uses only one
negative-odd spectral center, not an unspecified collection of
Stieltjes constants. The coefficient rule
\eqref{ag:eq:hrecurrence} makes the formula executable at every order.
No numerical independence assertion is needed for any of these identities.

% END sections/09_bernoulli.tex


% BEGIN sections/10_audit.tex
\section{Source audit, verification, and integration status}
\label{ag:sec:audit}

\subsection{What was inspected and what was not}
The GitHub metadata anchor used for the final source inventory is
\[
 \texttt{\pin}.
\]
The relevant canonical chapters, the immediate predecessor's Gamma
sector, and its research-question section were read as TeX through the
repository connection. The source manifest in the accompanying package
records exact paths and the returned Git blob identifiers. Some initial
reads used the default branch; the predecessor's research questions and
Gamma-section identity were subsequently checked at the displayed
anchor. A Git blob identifier is not being presented as a SHA-256 hash
of a locally downloaded file.

The current incoming directory contained five ZIP archives, concerning
alternating tails, finite-difference determinants, higher quadratic jets,
higher spectral defects, and quartic spectral reconstruction. The
repository connection exposed their metadata but not their binary
contents. The ZIPs were \emph{not} extracted in this continuation.
Matching-title Library TeX abstracts were retrieved for the
finite-difference, higher-quadratic-jet, and higher-spectral-defect
reports. Their identity with the repository ZIP bytes was not
established. These abstracts concern determinantal or spectral
factorization problems; they are contextual comparisons, not audited
proof dependencies. No assertion that every incoming result has been
read, or that the entire repository has been certified, is made.

The main proof dependencies are therefore unusually explicit: the
functional-monomial decomposition, the Gamma Bell exponential, elementary
weighted integration by parts, and the standard Hurwitz and Lerch
identities. The necessary algebraic decomposition is reproved here.
The one concrete predecessor question used to measure progress is
question~5 of \cite[Section~8]{HTG}, concerning the next Gamma
projection. Repository files have not been modified remotely.

\subsection{Classical inputs and the boundary of the novelty claim}
The functional-monomial construction is an established component of the
integro-differential-algebra literature \cite{GRR}; its adoption in the
canonical manuscript and predecessor is retained with attribution
\cite{Manuscript,HTG}. The pure derivative-square coefficients, the
low-index function $f_0$, and the entire $r=2$ family are inherited
specializations. They should not be relabeled as new theorems during
integration.

Polynomially weighted Hurwitz integrals and their relations to
Bernoulli polynomials, log Gamma, and polygamma functions have extensive
classical antecedents; the papers of Espinosa and Moll explicitly
study definite and indefinite integrals of this kind \cite{EM1,EM2}.
The differentiation, translation, Mellin, Bernoulli, and functional-
equation identities used here are standard \cite{DLMFzeta,DLMFpsi,
DLMFlerch,DLMFBernoulli,DLMFintegers}. Repeated integration by parts
itself is not a new method.

The proposed additions relative to the inspected proof chain are the
all-$r$ affine square-sector coefficient theorem, the arbitrary symmetric
profile extension, their common beta--Hurwitz transport, and the
explicit selector classification as a statement over $\C(s)$. The
all-order Stieltjes coefficient rules and anchored primitive formulas
are proved consequences in the stated normalization. Particular
low-order integral values may already occur in the literature.
This report does not establish global historical priority for each
formula, nor does it claim a new special-function algebraic-independence
result at fixed numerical arguments.

\subsection{The exact proof-status distinctions}
The following status list should accompany any manuscript integration.

\paragraph{Proved here.}
Theorem~\ref{ag:thm:main} gives every
$u_0^a u_1^k u_r^2$ coefficient in the canonical inverse-Gamma normal
form. Section~\ref{ag:sec:profiles} gives its general symmetric-profile
analogue. The Hurwitz and Lerch transforms, their removable-pole
prescriptions, all Stieltjes orders, and all anchored shift primitives
are analytic identities on their specified domains. Theorem~
\ref{ag:thm:selector-classification} classifies exactly the symmetric,
endpoint-vanishing polynomial weights with the specified single-zeta
functional closure. Its necessity uses a proof of independence over
$\C(s)$, rather than independence of a declared formal list.

\paragraph{A source question partially resolved.}
The next-Gamma-projection question is answered for $u_3^2$ with arbitrary
powers of $u_0,u_1$ and no additional $u_2$ factor, and simultaneously
for every higher derivative square. Terms such as
$u_0^a u_1^k u_2^b u_3^2$ with $b>0$ lie beyond the quadratic
high-derivative sector used here. The complete $\sigma_3$ projection is
not claimed.

\paragraph{Not resolved by this report.}
The unrestricted unequal-even-tail optimal-prefix conjecture, the
Gaussian $S_6$ target, and the current revised $S_8$ candidate remain
separate. The formal Gamma theorem does not evaluate all of its
finite-part moments. No proof-assistant formalization, independent
human peer review, or certified numerical enclosure is claimed.

\subsection{Corrections, safeguards, and editorial observations}
No false inherited mathematical theorem was established by the targeted
audit. Several errors would arise from tempting but invalid extensions,
and the package supplies explicit guards against them.
First, normal reduction must precede truncation of higher derivatives:
$\NF(u_0u_4)=u_2^2$, whereas truncating $u_4$ beforehand gives zero.
Second, $\NF$ is not multiplicative. Third, weighted symbol reduction
requires the Appell primitive property in
Lemma~\ref{ag:lem:appell-primitive}; allowing an arbitrary
$x$-dependent primitive would destroy the quotient argument. Fourth,
removable poles in endpoint formulas must be canceled jointly, not
by assigning independent finite parts to their summands. Finally,
$B_1(1)=+1/2$ and $B_1(0)=-1/2$ play different roles in the selector
calculation. These are proof obligations and implementation safeguards,
not accusations that the predecessor omitted them.

A minor editorial duplication was observed in the inspected manuscript
README: the paragraph describing the full-Cayley, 1,032-input first-$X$
family appeared twice. Removing one copy would improve the overview
without changing any mathematical claim. That suggestion should be
checked against the current README before editing, since the repository
is changing rapidly. No correction to a numbered theorem is proposed
on the basis of that duplication.

\subsection{Exact finite verification and numerical diagnostics}
The executable checks distinguish four kinds of exact finite assertion.
\begin{center}
\small
\begin{tabular}{@{}p{0.25\linewidth}rp{0.57\linewidth}@{}}
\toprule
Replay & Assertions & Scope\\
\midrule
\src{check_algebra.py}&1,255&Normal forms and primitives through weight 18;
independent Bell sectors; all-index formula specializations; kernel
coefficient identities; deliberately wrong alternatives rejected.\\[4pt]
\src{check_endpoint_algebra.py}&728&Beta endpoint derivatives and residues;
complete-harmonic coefficients; Bernoulli endpoint conventions and
rational offsets.\\[4pt]
\src{check_profiles.py}&75&Three independent polynomial profiles compared
with direct Bell expansion and canonical reduction.\\[4pt]
\src{check_selectors.py}&129&Selector basis, beta--Bernoulli transport,
Stieltjes residues, inverse expansions, and an endpoint-sign guard.\\
\midrule
Total&2,187&Finite exact rational or symbolic assertions.\\
\bottomrule
\end{tabular}
\end{center}
The algebra replay compares both sides with independently constructed
representations. In particular it checks $A_m=N_m+\D T_m$, not merely
the displayed selected coefficient. The profile replay does not invoke
the proposed profile formula to create its Bell-polynomial reference.
The infinite-family assertions nevertheless rest on the mathematical
proofs in the article, not the largest computed index.

The numerical program has four separately runnable parts: 15 kernel,
14 transform, 6 Stieltjes, and 21 Lerch/selector comparisons, totaling
56. At 50 working decimal digits it compares beta integrals with
hypergeometric values, compact Hurwitz averages with finite endpoint
expressions, Mellin integrals with both, and polygamma/Stieltjes
quadratures with their explicit evaluations. It includes a complex
spectral argument and rational-endpoint root-of-unity filters.
Errors are scaled by $\max(1,|L|,|R|)$; small absolute values therefore
do not imply the same number of relative digits.

The largest recorded scaled residual is approximately
$2.62\times10^{-29}$ in a comparison of two \emph{different quadrature
orders} for the first Stieltjes average. Its declared tolerance is
$10^{-23}$. The comparison of the finer quadrature with the finite
endpoint formula is about $5.49\times10^{-43}$, and the explicit
quartic $\gamma_1$ evaluation agrees to about $4.22\times10^{-40}$
in the same scaled measure. The default tolerance for the other
comparisons is $10^{-30}$, with explicit local overrides documented
in the program. These are floating-point diagnostics, not interval
error bounds.

All scripts, exact coefficient tables, normal-form primitives,
individual results, environment information, and a replay entry point
are included. The integration fragment preserves unique \src{ag:}
labels, the distinction between ordinary and finite-part integrals,
and the nonmultiplicative meaning of the normal projection.

% END sections/10_audit.tex


% BEGIN sections/11_questions.tex
\section{Further research questions and concrete next targets}
\label{ag:sec:questions}

The following problems continue the exact-identity programme rather
than asking for bounds or extrapolating a numerical pattern. The
proved formulas give initial conditions and explicit checks for each.

\paragraph{1. Complete the coupled next Gamma sector.}
Find a finite expression for the coefficient
\[
 [u_0^a u_1^k u_2^b u_3^2]N_{a+2k+3b+8}
\]
for every $b>0$.
The present theorem supplies $b=0$, while the predecessor supplies the
pure $u_2$ projection. A quadratic background for $u_0$ would retain
$u_2$ but replace constant-coefficient weighted adjoints by variable-
coefficient ones. The precise missing step is a canonical symbol
calculus compatible with that background and the original normal basis.
A formula should include the parity zeros and the retained endpoint
primitive, not just the first few nonzero weights.

\paragraph{2. Determine the first genuinely nonlinear polarization.}
The affine quadratic symbol is a polynomial in one invariant $\eta$.
At cubic and quartic perturbation order there are several symmetric
invariants after imposing $\sum\xi_i+\kappa=0$. Identify an invariant
basis which maps triangularly to the functional-monomial normal form.
Use it to compute complete families with three or four high-derivative
factors, distinguishing cubic terms such as $u_2^3$ from quartic terms
such as $u_2^2u_3^2$. The source's all-$u_2$ formulas provide stringent
independent checks but not the mixed-sector solution.

\paragraph{3. Classify elementary symmetric profiles.}
For which even analytic profiles $K$ do all square-sector kernels
$g_{K,r}$ belong to a fixed finite combination of hyperbolic functions
and their divided differences? The finite Gamma mixtures in
Section~\ref{ag:sec:profiles} give one class. Seek a converse under
explicit growth and analyticity hypotheses, or exhibit another class
coming from an interpretable Gamma or Barnes ratio. A claim about
minimality must specify the allowed coefficient field and operations.

\paragraph{4. Generalize the selector classification to other intervals.}
On $[1,3]$, integer translation leaves only the bases $1$ and $2^{-s}$
besides shifted Riemann zeta functions. For intervals $[A,B]$ with
positive integer endpoints, prescribe the finite exponential remainder
and characterize weights selecting exactly one shifted zeta function.
For rational endpoints the natural coordinates are Hurwitz values or
Dirichlet $L$-functions. The goal is an exact functional classification,
including endpoint multiplicities and a proof of the necessary
Dirichlet-series independence statement.

\paragraph{5. Find Lerch selectors with a prescribed polylogarithm order.}
The general Lerch endpoint formula is finite, but the rational factors
$z^{-A}$ prevent the simple odd-derivative diagonalization used at $z=1$.
Characterize polynomial weights, possibly depending rationally on $z$,
whose transform has one polylogarithm order plus an elementary
remainder. Determine whether the required endpoint conditions form a
finite interpolation problem, and prove its confluence as $z\to1$.
The explicit quartic identity in Section~\ref{ag:sec:polylogs} is a
normalization test rather than a general classification.

\paragraph{6. Put all selector Stieltjes orders in a uniform closed form.}
Equation~\eqref{ag:eq:selector-Stieltjes} is a finite coefficient rule
at every order. Seek a compact bivariate generating identity in the
selector degree and Stieltjes order that preserves the single spectral
center and makes the harmonic coefficients explicit. Determine which
anchored integrals in the shift can be expressed through a fixed family
of multiple-Gamma or zeta-derivative functions. Any such primitive must
retain the anchor conditions from Section~\ref{ag:sec:stieltjes}.

\paragraph{7. Describe the numerical image of the Gamma square sector.}
The coefficient theorem is formal. Evaluate or systematically reduce
the ordinary shifted integrals and the original-coordinate finite parts
of $\psi^a(\psi')^{2k}(\psi^{(r)})^2$ in a declared basis of convergent
multiple-zeta derivatives or polylogarithmic integrals. The derivative-
ideal pole calculus in the manuscript supplies convergent coordinates,
but it does not establish their arithmetic minimality. Retain the
primitive endpoint term when transporting the result to inverse-Gamma
moments.

\paragraph{8. Return to the unequal-tail boundary determinant.}
The all-$r$ Gamma theorem does not prove the unequal-even-tail
optimal-prefix conjecture. The next unresolved general fixed-gap
problem in the inspected predecessor is the three-column boundary
determinant for gap six. Find an exact scalar recurrence, a determinant
identity, or an arithmetic normalization that proves nonvanishing.
Entrywise positivity and pairwise ratio separation are not sufficient.
Any resulting inverse row can then be combined with the predecessor's
already proved universal finite telescoper.

\paragraph{9. Build a proof-carrying coefficient compiler.}
Formalize the terminating normal-form algorithm and its identity
$P=\NF(P)+\D\Prim(P)$ in a proof assistant. Add the Appell covariance
lemma and the weighted quadratic symbol as reusable algebraic results.
A useful initial target is a checker for the supplied rational JSON
certificates, followed by the all-order coefficient proof. Certified
real or complex interval checks of the displayed analytic integrals
would be a separate numerical project, not a replacement for analytic
continuation proofs.

\paragraph{10. Preserve the original period questions.}
The Gaussian $S_6$ and current revised $S_8$ candidates call for exact
polylogarithmic identities with fixed branches and complete lower-weight
constants. Neither selector closure nor normal-form coefficient
closure settles them. A promising next step is a parameter identity
whose rigorously controlled boundary value matches the precise frozen
candidate, rather than a new unconstrained integer-relation search.

\subsection*{Conclusion}
An affine perturbation converts an infinite family of Gamma normal-form
coefficients into one entire divided difference. The beta representation
of its coefficients supplies a common route to finite Hurwitz and Lerch
endpoint formulas, generalized Stieltjes averages, and anchored
primitives. Bernoulli endpoint data then diagonalize the integer-
interval arithmetic and give a genuine single-zeta functional
classification. These are exact identities with separate formal,
analytic, and computational meanings. Their limitations identify the
next nonlinear and arithmetic questions without weakening the proved
all-index statements.

% END sections/11_questions.tex


% BEGIN references.tex
\begin{thebibliography}{99}
\raggedright
\bibitem{Manuscript}
ProveIt Contributors, \emph{Polylogarithms and their Arithmetic Bridges}.
Canonical manuscript and the chapters
\src{07-ibp-normal-forms.tex} and
\src{07-differential-polynomial-moments.tex}.
Repository metadata anchor \texttt{\pin}; consulted 11 October 2026.
\url{https://github.com/VladimirReshetnikov/ProveIt/tree/92ee3aaf3087c0e0c29b9ed952e777de5d6accd0/Analysis/Polylogarithms/docs/manuscript}.

\bibitem{HTG}
\emph{Exact Harmonic-Tail Reconstruction and Variational Gamma
Identities}. ProveIt working report, 11 October 2026.
Especially Section~4, \src{sections/04_gamma_sectors.tex}, and
Section~8, \src{sections/08_audit_research.tex}, research question~5.
\url{https://github.com/VladimirReshetnikov/ProveIt/tree/92ee3aaf3087c0e0c29b9ed952e777de5d6accd0/Analysis/Polylogarithms/docs/reports/stieltjes-harmonic-resonance/continuations/harmonic-tail-gamma}.
The report's internal historical baseline is different from this
continuation's metadata anchor; the source manifest records that distinction.

\bibitem{GRR}
Li Guo, Georg Regensburger, and Markus Rosenkranz,
\emph{On integro-differential algebras}.
Journal of Pure and Applied Algebra \textbf{218} (2014), no.~3,
456--473. \url{https://doi.org/10.1016/j.jpaa.2013.06.015}.
Preprint: \url{https://arxiv.org/abs/1212.0266}.
The functional-monomial construction is a classical input.

\bibitem{EM1}
Olivier R. Espinosa and Victor H. Moll,
\emph{On some definite integrals involving the Hurwitz zeta function}.
Preprint, 2000. \url{https://arxiv.org/abs/math/0012078}.
Cited for the established literature on Hurwitz, Bernoulli, and
log-Gamma integrals; not as a source of the new affine coefficient theorem.

\bibitem{EM2}
Olivier R. Espinosa and Victor H. Moll,
\emph{On some integrals involving the Hurwitz zeta function: part 2}.
Preprint, 2001. \url{https://arxiv.org/abs/math/0107082}.
Cited for the established indefinite-integral and negative-polygamma
programme.

\bibitem{DLMFpsi}
NIST Digital Library of Mathematical Functions, \S5.15,
\emph{Polygamma Functions}. Consulted 11 October 2026.
\url{https://dlmf.nist.gov/5.15}.

\bibitem{DLMFzeta}
NIST Digital Library of Mathematical Functions, \S25.11,
\emph{Hurwitz Zeta Function}. Consulted 11 October 2026.
Definitions, Mellin representation, argument differentiation,
integer translation, Bernoulli values, and the log-Gamma specialization.
\url{https://dlmf.nist.gov/25.11}.

\bibitem{DLMFintegers}
NIST Digital Library of Mathematical Functions, \S25.6,
\emph{Integer Arguments}. Consulted 11 October 2026.
Zeta values and derivatives at nonpositive integers.
\url{https://dlmf.nist.gov/25.6}.

\bibitem{DLMFlerch}
NIST Digital Library of Mathematical Functions, \S25.14,
\emph{Lerch's Transcendent}. Consulted 11 October 2026.
Defining series, Mellin representation, and polylogarithmic specialization.
\url{https://dlmf.nist.gov/25.14}.

\bibitem{DLMFBernoulli}
NIST Digital Library of Mathematical Functions, \S24.2,
\emph{Definitions and Generating Functions}, and \S24.4,
\emph{Basic Properties}. Consulted 11 October 2026.
\url{https://dlmf.nist.gov/24.2}; \url{https://dlmf.nist.gov/24.4}.
\end{thebibliography}

% END references.tex

\end{document}
